\chapter{Measures on Hausdorff topological spaces}

If T is a set, and A is a subset of T, we denote by \(\varphi_{A}\) the characteristic function of A, provided this does not lead to any confusion. The set \(\overline{R}_{+}^{T}\) of numerical functions \(\geqslant0\) (finite or not) defined on T will be denoted by \(\mathcal{F}_{+}(T)\) , or simply \(F_{+}\) if there is no ambiguity as to T; this set will always be equipped with its natural order structure. Recall that the product of two elements of \(F_{+}\) is always defined, thanks to the convention \(0\cdot(+\infty)=0\) . If A is a subset of T, and f is a function defined on T, the restriction \(f|_{A}\) of f to A may be denoted \(f_{A}\) in this chapter, if this creates no confusion; an analogous notation will be employed for induced measures. On the other hand, if \(f\in\mathcal{F}_{+}(A)\) we shall denote by \(f^{0}\) the extension by 0 of f to T, that is, the function defined on T that coincides with f on A and with 0 on T-A.

All topological spaces considered in this chapter are assumed to be Hausdorff, absent express mention to the contrary. From §1, No.~4 on, except for §5, all measures will be assumed to be positive, absent express mention to the contrary.

